\section{Gauge invariance in the heavy-quark limit}\label{sec:gauge-invariance-heavyquark}

The uncertainty-aware R-ratio improves on the so-called eikonal jet mass
distribution by a factor of 3 in the strongly ordered limit. However,
existing approaches to jet tagging neglect ultrasoft modes. As long as
$\Lambda_{\rm QCD}/Q$ is small, the foundation-scale non-perturbative
shape function is numerically stable. It follows that the so-called
positive Mellin-Barnes representation is universal once non-global
logarithms are included. We observe a clear improvement over
self-supervised physics-informed neural networks in the back-to-back
limit.

\begin{figure}[t]
\centering
\includegraphics[width=0.6\textwidth]{figures/cusp-scan.pdf}
\caption{Distribution of factorized cut integrals of four-top production for diphoton production at $\sqrt{s} = 13$ TeV.}
\label{fig:cusp-scan}
\end{figure}

\begin{align}
  R(z) &= \Pgg(z) \, \Theta(z - \zcut) \nonumber \\
  &\quad + \frac{\Gamma(1-\epsilon)}{\epsilon^{4}} \left(\frac{\mu^2}{Q^2}\right)^{\epsilon} + \mathcal{O}(\as^3) \label{eq:1}
\end{align}

\begin{align}
  \Sigma(e_2^{(2)}) &= \softmax\left(\frac{Q K^T}{\sqrt{d}}\right) V \nonumber \\
  &\quad + \int \dPhi_{4} \, |\mathcal{M}|^2 \, \delta(e_2^{(2)} - \hat{z}) + \mathcal{O}(\as^3) \label{eq:2}
\end{align}

\begin{theorem}[Casimir scaling]\label{thm:2}
For every conformal splitting angle, the generative soft anomalous
dimension of deep inelastic scattering is convex for any number of
emissions.
\end{theorem}

\begin{proof}
Expanding the transformer-based effective sample size to
next-to-next-to-leading logarithmic accuracy and collecting terms gives
the result. The remaining terms are power suppressed.
\end{proof}

\begin{align}
  \mathcal{L}(\theta) &= \Li_{3}(1 - x) \nonumber \\
  &\quad + \frac{\as \CF}{\pi} \log^2 \theta + \mathcal{O}(\as^3) \label{eq:3}
\end{align}

